% META: {"id": "1NSI-LANG-EVAL-B", "chapitre": "1NSI-LANGAGE", "type_objet": "evaluation", "capacites": ["P-LANG-01", "P-LANG-02", "P-LANG-03A", "P-LANG-03B", "P-LANG-03C", "P-LANG-04", "P-LANG-05"], "mode_creation": "ex_nihilo", "sources_inspiration": [], "duree_min": 60, "status": "verified"}
\section*{Évaluation B --- Langages et programmation}
\textit{Durée : 45 minutes. Ordinateur avec interpréteur Python autorisé.}

\baremeIndicatif{Ex. 1 : 10 pts (constructions, spécification) ; Ex. 2 : 10 pts (mise au point, bibliothèque) ; Ex. 3 : 10 pts (traits communs et particuliers, prototype, postcondition)}

\bigskip

\textbf{Exercice 1} --- \textit{Moyenne pondérée} \hfill (10 points)

\begin{enumerate}
  \item (2 pts) Proposer une précondition pour une fonction \lstinline{moyenne_ponderee(valeurs, poids)} calculant une moyenne pondérée.
  \item (5 pts) Écrire cette fonction, avec les \lstinline{assert} correspondants, en utilisant une boucle \lstinline{for} bornée.
  \item (3 pts) Calculer \lstinline{moyenne_ponderee([8, 12, 16], [1, 1, 2])} à la main, puis vérifier avec le code.
\end{enumerate}

% BEGIN-VERIFY
% def moyenne_ponderee(valeurs, poids):
%     assert len(valeurs) == len(poids)
%     assert sum(poids) > 0
%     total = 0
%     for i in range(len(valeurs)):
%         total = total + valeurs[i] * poids[i]
%     return total / sum(poids)
% assert moyenne_ponderee([8, 12, 16], [1, 1, 2]) == 13.0
% END-VERIFY

\bigskip

\textbf{Exercice 2} --- \textit{Mise au point et bibliothèque} \hfill (10 points)

Un élève écrit la fonction suivante, censée renvoyer le produit des éléments strictement
négatifs d'une liste :
\begin{python}
def produit_negatifs_bugue(liste):
    prod = 0
    for x in liste:
        if x < 0:
            prod *= x
    return prod
\end{python}

\begin{enumerate}
  \item (3 pts) Que renvoie \lstinline{produit_negatifs_bugue([-2, -3, -4])} ? Quel est le résultat attendu ?
  \item (4 pts) Identifier précisément l'erreur dans le code, et proposer une version corrigée.
  \item (3 pts) En utilisant le module \lstinline{math}, calculer le PGCD de $60$ et $48$ (indiquer le nom de la fonction utilisée).
\end{enumerate}

% BEGIN-VERIFY
% def produit_negatifs_bugue(liste):
%     prod = 0
%     for x in liste:
%         if x < 0:
%             prod *= x
%     return prod
% assert produit_negatifs_bugue([-2, -3, -4]) == 0
% def produit_negatifs(liste):
%     prod = 1
%     for x in liste:
%         if x < 0:
%             prod *= x
%     return prod
% assert produit_negatifs([-2, -3, -4]) == -24
% import math
% assert math.gcd(60, 48) == 12
% END-VERIFY

\bigskip

\textbf{Exercice 3} --- \textit{D'un langage à l'autre, et le contrat d'une fonction} \hfill (10 points)

On donne cet extrait, écrit dans un langage à accolades :

\begin{console}
function compte(t, x) {
    let n = 0;
    for (let i = 0; i < t.length; i = i + 1) {
        if (t[i] === x) { n = n + 1; }
    }
    return n;
}
\end{console}

\begin{enumerate}
  \item (3 pts) Citer trois éléments du \textbf{noyau commun} présents dans cet extrait, et
        trois éléments \textbf{propres} à ce langage.
  \item (3 pts) Réécrire cette fonction en Python, sans changer son comportement.
  \item (2 pts) Écrire le \textbf{prototype} de cette fonction : nom, paramètres, valeur
        renvoyée. Ne rien dire de la manière dont elle procède.
  \item (2 pts) Énoncer une \textbf{postcondition} portant sur le résultat, puis expliquer
        pourquoi « les éléments de \lstinline{t} doivent être comparables à \lstinline{x} »
        n'en est pas une.
\end{enumerate}

% BEGIN-VERIFY
% def compte(t, x):
%     n = 0
%     for v in t:
%         if v == x:
%             n = n + 1
%     return n
% t = [3, 7, 3, 3, 9]
% r = compte(t, 3)
% assert r == 3
% assert 0 <= r <= len(t)
% assert compte(t, 42) == 0
% assert compte([], 1) == 0
% assert t == [3, 7, 3, 3, 9]
% END-VERIFY
